\clearpage\thispagestyle{empty}
\pdfbookmark[0]{Quellen zu Pallas (Brief vom 15. April 1812)}{pallas1812apr15-sources}
\begin{center}\textbf{Quellen zu Pallas (Brief vom 15. April 1812)}\end{center}
\noindent Diese KI-Übertragung ins heutige Deutsch folgt den \href{https://gdz.sub.uni-goettingen.de/id/PPN236007467}{\textit{Werken}, Band~VI} (Göttingen, 1874), S.~349. Gauß’ Brief vom 15. April erschien als Nr.~XXXVI, \textit{Aus einem Schreiben des Hrn. Prof. Gauss}, in der \textit{Monatlichen Correspondenz}, Band~XXV, April 1812, S.~\href{https://api.digitale-sammlungen.de/iiif/image/v2/bsb10538614_00403/full/1600,/0/default.jpg}{389}. Das digitalisierte Exemplar stammt aus der Bayerischen Staatsbibliothek.
\par\smallskip\noindent Die Schreibweise folgt der Werkausgabe. In der Beobachtungszeile steht das Zeichen für Pallas. Der Ausdruck im Absatz über die Störungen enthält die Zeichen für Jupiter und Pallas. KI-Übertragung und KI-Satz.
\clearpage
